\documentclass[10pt]{article}

% ---------- Document Identity (update these only) ----------
\newcommand{\DocStage}{}
\newcommand{\DocTitle}{Your Road to Tier One SOF}
\newcommand{\ProgramName}{182-Week Civilian-to-Selection Readiness Pipeline}
\newcommand{\ProgramSubtitle}{}

% ---------- Page ----------
\usepackage[legalpaper,left=0.5in,right=0.5in,top=0.6in,bottom=0.45in]{geometry}

% ---------- Typography ----------
\usepackage[T1]{fontenc}
\usepackage[utf8]{inputenc}
\usepackage{libertinus}
\usepackage{microtype}
\usepackage{parskip}
\usepackage{ragged2e}
\usepackage{textcomp}
\AtBeginDocument{\color{black}}
\AtBeginDocument{\pagecolor{white}}
\AtBeginDocument{\justifying}

% ---------- Landscape pages ----------
\usepackage{pdflscape}

% ---------- Color ----------
\usepackage[table]{xcolor}
\definecolor{StageBlue}{HTML}{1F4E79}
\definecolor{StageBlueDark}{HTML}{163A5A}
\definecolor{LightGray}{HTML}{F2F4F7}
\definecolor{MidGray}{HTML}{D9DEE5}
\definecolor{RuleGray}{HTML}{C7CED8}
\definecolor{HeaderInk}{HTML}{000000}
\definecolor{Ink}{HTML}{000000}

% ---------- Utility ----------
\usepackage{enumitem}
\sloppy
\setlength{\emergencystretch}{2em}

% ---------- Boxes / UI ----------
\usepackage[most]{tcolorbox}

\tcbset{
  charterTitle/.style={
    enhanced,
    colback=StageBlue!4,
    colframe=StageBlue,
    boxrule=1.25pt,
    arc=3pt,
    left=16pt,right=16pt,top=14pt,bottom=14pt,
    borderline north={3.2pt}{0pt}{StageBlue},
    borderline west={4pt}{0pt}{StageBlue},
    borderline south={1.25pt}{0pt}{StageBlue},
    drop shadow={black!25!white},
  },
  charterRule/.style={
    enhanced,
    colback=LightGray,
    colframe=RuleGray,
    boxrule=0.85pt,
    arc=3pt,
    left=12pt,right=12pt,top=10pt,bottom=10pt,
    borderline west={3pt}{0pt}{StageBlue},
  },
  charterBadge/.style={
    enhanced,
    colback=StageBlue,
    coltext=white,
    colframe=StageBlueDark,
    boxrule=0.6pt,
    arc=2pt,
    left=6pt,right=6pt,top=3pt,bottom=3pt,
  }
}
\newtcbox{\Badge}{nobeforeafter,charterBadge}

% ---------- Lists ----------
\setlist[itemize]{leftmargin=1.5em, itemsep=0.25em, topsep=0.25em}
\setlist[enumerate]{leftmargin=1.8em, itemsep=0.25em, topsep=0.25em}

% ---------- TOC ----------
\usepackage{tocloft}
\setcounter{tocdepth}{3}
\setlength{\cftbeforesecskip}{3pt}
\setlength{\cftbeforesubsecskip}{2pt}
\setlength{\cftbeforesubsubsecskip}{0pt}
\renewcommand{\cftsecfont}{\bfseries}
\renewcommand{\cftsecpagefont}{\bfseries}
\renewcommand{\cftsubsecfont}{\normalfont}
\renewcommand{\cftsubsecpagefont}{\normalfont}
\renewcommand{\cftsubsubsecfont}{\normalfont}
\renewcommand{\cftsubsubsecpagefont}{\normalfont}
\setlength{\cftsecindent}{0pt}
\setlength{\cftsubsecindent}{1.25em}
\setlength{\cftsubsubsecindent}{2.8em}
\setlength{\cftsecnumwidth}{2.1em}
\setlength{\cftsubsecnumwidth}{2.8em}
\setlength{\cftsubsubsecnumwidth}{3.6em}

% Remove the built-in TOC heading so only the custom heading is shown
\makeatletter
\renewcommand{\tableofcontents}{\@starttoc{toc}}
\makeatother

% ---------- Header/Footer: page number only ----------
\usepackage{fancyhdr}
\pagestyle{fancy}
\fancyhf{}
\renewcommand{\headrulewidth}{0pt}
\renewcommand{\footrulewidth}{0pt}
\fancyfoot[C]{\thepage}

% ---------- Section formatting ----------
\usepackage{titlesec}

\titleformat{\section}
  {\Large\bfseries\color{Ink}}
  {\thesection}{0.6em}{}
  [\vspace{0.25em}\textcolor{StageBlue}{\rule{\linewidth}{0.8pt}}]

\titleformat{\subsection}
  {\large\bfseries\color{Ink}}
  {\thesubsection}{0.6em}{}

\titleformat{\subsubsection}
  {\normalsize\bfseries\color{Ink}}
  {\thesubsubsection}{0.6em}{}

\titlespacing*{\section}{0pt}{0.95em}{0.35em}
\titlespacing*{\subsection}{0pt}{0.65em}{0.25em}
\titlespacing*{\subsubsection}{0pt}{0.55em}{0.2em}

% ---------- Table engine ----------
\usepackage{tabularray}
\UseTblrLibrary{booktabs}

% ---------- tabularray: GLOBAL table treatment ----------
\SetTblrInner{
  cells = {fg=black, bg=white, valign=t},
}

% ---------- Header helpers ----------
\newcommand{\Hdr}[1]{\shortstack[c]{\strut #1\strut}}

% ---------- Lines ----------
\newcommand{\TitleLine}{\textcolor{StageBlue}{\rule{0.98\linewidth}{0.95pt}}}
\newcommand{\SubTitleLine}{\textcolor{RuleGray}{\rule{0.98\linewidth}{0.65pt}}}

% ---------- Continued on next page ----------
\newcommand{\ContinuedOnNextPageInline}{%
  \par
  \vfill
  \noindent\textcolor{RuleGray}{\rule{\linewidth}{1pt}}\vspace{-2pt}
  \noindent\hfill\textit{Continued on next page}\par\vspace{-10pt}
  \noindent\textcolor{RuleGray}{\rule{\linewidth}{1pt}}
  \clearpage
}

% ---------- PDF hyperlinks ----------
\usepackage[hidelinks]{hyperref}
\hypersetup{
  colorlinks=false,
  pdfborder={0 0 0},
  linktoc=all,
  pdftitle={\DocTitle},
  pdfauthor={\ProgramName}
}

% =========================================================
\begin{document}

\section*{Stage 3.E — Adjustment Rules for Illness, Regression, and Risk Flags}
\phantomsection
\addcontentsline{toc}{section}{Stage 3.E — Adjustment Rules for Illness, Regression, and Risk Flags}

\subsection*{3.E.1 Core Adjustment Principles}
\phantomsection
\addcontentsline{toc}{subsection}{3.E.1 Core Adjustment Principles}

\begin{itemize}
\item Validity-first is mandatory: \textbf{INVALID} tests provide no gate credit and cannot be used to justify \textbf{ADVANCE}; the corrective action is reslot to the next protected deload+test window.
\item Safety-first is mandatory: red-flag symptoms or safety failures trigger \textbf{STOP} \textbf{NOW}; the day’s planned tests are aborted and are inadmissible for gate use; \textbf{MEDICAL} \textbf{HOLD} posture is applied when indicated by stop-rule triggers or clinician restriction.
\item Preserve macro integrity: keep the work-to-deload rhythm intact; adjust by reslotting tests rather than compressing, stacking, or forcing tests into build weeks.
\item No stacking make-ups: missed or invalid tests are not “recovered” by stacking multiple high-cost tests into one week or by violating Stage 3.A spacing intent.
\item Reslot priority order (fixed; use in this sequence):
\begin{enumerate}
\item next deload+test week within the same phase,
\item if the end gate is missed: use an existing buffer or Transition week if explicitly defined in Stage 3.D; otherwise extend the phase under patch control,
\item if still not possible or risk posture remains unacceptable: default to \textbf{HOLD}, or \textbf{REGRESS}, or \textbf{MEDICAL} \textbf{HOLD} consistent with Stage 1 governance.
\end{enumerate}
\item Environment validity is a gate control: unsafe or non-standard route/pool/machine conditions trigger postpone and reslot; do not accept compromised conditions as “good enough” for gate-grade evidence.
\item Tool drift is validity drift: undocumented tool changes (route verification method, pool unit/length, treadmill identity/settings, timing method) invalidate gate-grade interpretation and require either verified late-entry proof or reslot.
\item First-appearance posture is protected: higher-cost proof events are not pulled forward into earlier segments as a substitute after a miss; they are reslotted, not escalated.
\item Convert, do not force: when readiness is compromised, convert a planned deload+test week into deload-only or check-in-only; reslot maximal tests forward.
\item Retest discipline is mandatory: retests occur only in deload+test weeks or reduced-stress windows and must respect Stage 3.A stacking and spacing intent.
\item Documentation is a control: every adjustment \textbf{MUST} be logged with (a) validity tag impact (\textbf{VALID}/\textbf{MODIFIED}/\textbf{INVALID} for sessions; \textbf{VALID}/\textbf{INVALID} for tests), (b) reason-code posture (\textbf{SESSION-RC} or \textbf{TEST-RC} as applicable), (c) gate impact posture (\textbf{ADVANCE}/\textbf{HOLD}/\textbf{REGRESS}/\textbf{MEDICAL} \textbf{HOLD}), and (d) reslot target week.
\item Domain cost hierarchy for overload resolution:
\begin{itemize}
\item protect the primary proof event first (the decisive gate item for the phase),
\item defer secondary high-cost events before deferring the primary proof event,
\item maintain continuous low-cost monitoring (\textbf{BODYCOMP}, \textbf{RECOVERY}, \textbf{DURABILITY}) even when major tests are deferred.
\end{itemize}
\item Underwater governance removal rule (fixed): removing underwater from a scheduled window because governance/supervision is unavailable is allowed without patch; it \textbf{MUST} be documented as “not scheduled; inadmissible without governance; surface-only schedule continues,” and no proxy equivalency is permitted.
\end{itemize}

\subsection*{3.E.2 Illness and Acute Injury Handling Rules}
\phantomsection
\addcontentsline{toc}{subsection}{3.E.2 Illness and Acute Injury Handling Rules}

\begin{itemize}
\item Cancel/postpone criteria (non-diagnostic; governance-level):
\begin{itemize}
\item \textbf{STOP} \textbf{NOW} triggers (chest pain/pressure; syncope/near-syncope; disproportionate shortness of breath; new neurologic symptoms; heat illness signs; acute injury with instability/swelling; gait-altering pain/mechanics change) cancel all testing that day and render intended tests inadmissible for gate use.
\item Illness with systemic symptoms (e.g., feverishness, severe GI distress, significant respiratory illness reducing exertion tolerance) postpones gate-grade testing and converts the planned test week to deload-only or check-in-only.
\end{itemize}

\item What gets logged (minimum required):
\begin{itemize}
\item Test Validity Tag outcome: planned gate-grade test is recorded as not executed or aborted; any partial attempt is \textbf{INVALID} for gate use.
\item Reason-code posture:
\begin{itemize}
\item Sessions: use \textbf{SESSION-RC-003} (\textbf{ILL}) or \textbf{SESSION-RC-002} (PAIN) or \textbf{SESSION-RC-007} (\textbf{SAFETY}) as primary as applicable; add free-text trigger detail.
\item Tests: use \textbf{TEST-RC-004} (\textbf{SAFETY}) when aborted for safety; use \textbf{TEST-RC-002} (\textbf{ENV}) when environment is unsafe; use \textbf{TEST-RC-001} (\textbf{MEAS}) when measurement integrity fails.
\end{itemize}
\item Reslot target window: the next deload+test week within the same phase is the default.
\item Gate impact posture: default \textbf{HOLD} if end-of-phase gate cannot be executed validly.
\end{itemize}

\item Return-to-testing posture:
\begin{itemize}
\item Testing is delayed until symptoms resolve and no red-flag posture exists.
\item Near-max or maximal testing is delayed until safe; do not attempt “modified maximal” compromises to salvage a gate.
\end{itemize}

\item Multi-week disruption posture:
\begin{itemize}
\item Default gate posture is \textbf{HOLD}.
\item If the end gate cannot be executed within remaining protected windows, phase extension is permitted only under patch control.
\item If risk posture worsens or stop-rule events recur, gate posture may escalate to \textbf{REGRESS} or \textbf{MEDICAL} \textbf{HOLD} per governance.
\end{itemize}

\item Binding rule: red-flag symptoms on test day mean abort; no modified maximal test is attempted “to get something on the board.”
\end{itemize}

\subsection*{3.E.3 Regression Handling Rules}
\phantomsection
\addcontentsline{toc}{subsection}{3.E.3 Regression Handling Rules}

\begin{itemize}
\item Noise versus true regression (operational, conservative):
\begin{itemize}
\item Noise is a single off-day output with no corroborating pattern and no safety compromise; it is recorded but does not force immediate schedule compression.
\item Meaningful regression is any of the following:
\begin{itemize}
\item a \textbf{VALID} test result below the required tier threshold for the current gate context (do not redefine thresholds here),
\item two consecutive \textbf{VALID} results trending worse in a gate-critical metric within the decision window,
\item a \textbf{VALID} result accompanied by worsening \textbf{DURABILITY}/\textbf{RECOVERY} signals that indicate the performance is not safely reproducible.
\end{itemize}
\end{itemize}

\item Retest attempt posture (binding):
\begin{itemize}
\item There is no fixed cap on attempts because advancement requires standards; however, retests \textbf{MUST} remain controlled:
\begin{itemize}
\item \textbf{VALID} under comparable conditions and fully logged per Stage 2 conventions,
\item scheduled only in deload+test weeks, baseline weeks, or reduced-stress windows explicitly intended for verification,
\item must not violate Stage 3.A stacking and spacing intent,
\item repeated retests must not be compressed into consecutive weeks unless replacing an \textbf{INVALID}/missed attempt under controlled conditions and without increasing risk exposure.
\end{itemize}
\end{itemize}

\item What not to do (explicit prohibitions):
\begin{itemize}
\item Do not chase immediate retests repeatedly (“retest spiral”).
\item Do not stack multiple time trials across domains in the same week to compensate for one regression.
\item Do not generate a cycle of \textbf{INVALID} attempts to manufacture progress; \textbf{INVALID} provides no gate credit.
\end{itemize}

\item Gate posture linkage:
\begin{itemize}
\item If regression occurs on a gate-critical metric and cannot be corrected with a valid retest within the phase’s remaining protected windows, default posture is \textbf{HOLD} with reslot or extension under patch control.
\item If regression coincides with durability breakdown (gait-altering pain, recurrent foot breakdown, repeated \textbf{STOP} \textbf{NOW} triggers), escalation to \textbf{REGRESS} or \textbf{MEDICAL} \textbf{HOLD} is indicated per governance.
\end{itemize}
\end{itemize}

\subsection*{3.E.4 Risk Flags and Readiness-Based Adjustments}
\phantomsection
\addcontentsline{toc}{subsection}{3.E.4 Risk Flags and Readiness-Based Adjustments}

Hard-stop triggers (cancel/postpone; may initiate \textbf{MEDICAL} \textbf{HOLD}):

\begin{itemize}
\item \textbf{STOP} \textbf{NOW} red-flag symptoms (cardiovascular, neurologic, heat illness signs).
\item Gait-altering pain or mechanics change (\textbf{DURABILITY} constraint with immediate termination posture).
\item Acute instability or swelling suggesting acute injury risk.
\item Water governance violation (any unsupervised underwater/breath-hold attempt) or missing governance prerequisites for scheduled underwater work.
\item Unsafe environment for the planned test: ice/unsafe footing, severe wind/visibility hazards, lightning/storm risk, unsafe heat exposure, poor air quality, pool closure/crowding that breaks validity, machine failure or unverified settings.
\end{itemize}

Yellow-flag trend triggers (downgrade test week to deload-only or check-in-only; reslot maximal tests forward):

\begin{itemize}
\item Sustained sleep debt trend (persistent low sleep opportunity or worsening sleep quality) that increases injury risk and contaminates test validity.
\item Elevated soreness/fatigue trend that indicates fatigue contamination risk for gate-grade tests.
\item Worsening pain trend even if non-emergent, especially if it approaches mechanics change.
\item Repeated \textbf{MODIFIED} or \textbf{INVALID} sessions in the decision window (evidence-quality failure and/or recovery collapse).
\item Repeated missed sessions logged under time/environment constraints (\textbf{SESSION-RC-005} \textbf{TIME}; \textbf{SESSION-RC-001} \textbf{ENV}; \textbf{SESSION-RC-008} \textbf{WEATHER}) and/or \textbf{TRAVEL} disruptions.
\item Tool/environment drift not stabilized (new route, new treadmill, new pool unit/length, new footwear interface) inside a test window.
\end{itemize}

Governance-level examples (non-diagnostic; scheduling actions):

\begin{itemize}
\item If a deload+test week begins with yellow-flag trends (sleep collapse, rising soreness, repeated \textbf{MODIFIED} sessions), convert the week to deload-only or check-in-only:
\begin{itemize}
\item permit low-cost monitoring and one low-cost check-in major event only if it remains within Stage 3.A caps/spacing intent,
\item reslot maximal \textbf{RUN}/\textbf{RUCK}/\textbf{SWIM} time trials to the next protected window.
\end{itemize}
\item If a hard-stop trigger occurs during a scheduled test, terminate immediately; record the test as \textbf{INVALID} with \textbf{TEST-RC-004} (\textbf{SAFETY}) and default gate posture \textbf{HOLD} for that metric until a \textbf{VALID} retest exists.
\end{itemize}

\subsection*{3.E.5 Reslotting and Phase Extension Rules}
\phantomsection
\addcontentsline{toc}{subsection}{3.E.5 Reslotting and Phase Extension Rules}

Reslotting rules (binding):

\begin{itemize}
\item Missed or \textbf{INVALID} end-of-phase gate tests default to \textbf{HOLD}.
\item Required gate tests \textbf{MUST} be reslotted to the next available protected window inside the same phase.
\item If no protected window remains, phase extension is permitted only under patch control. There is no fixed maximum extension cap; extensions must remain conservative, validity-driven, and must not violate Stage 3.A stacking/spacing intent.
\item Underwater removal due to missing governance is allowed without patch and must be documented as not scheduled and inadmissible without governance; no proxy replacement.
\end{itemize}

Allowed without patch (must be documented and auditable):

\begin{itemize}
\item Swap sequencing of tests within the same deload+test week (without adding events, without violating caps/spacing intent).
\item Reslot a missed test to the next scheduled deload+test week within the same phase.
\item Downgrade a planned deload+test week into deload-only or check-in-only and reslot maximal tests forward.
\item Remove underwater testing when governance is unavailable, documented as inadmissible without governance; surface-only schedule continues.
\item Handle travel disruption using existing time or environment reason codes plus free-text note \textbf{TRAVEL}:
\begin{itemize}
\item Primary session cause \textbf{SHOULD} be \textbf{SESSION-RC-006} (\textbf{TRAVEL}) where applicable; add \textbf{SESSION-RC-005} (\textbf{TIME}) or \textbf{SESSION-RC-001} (\textbf{ENV}) only as explanatory detail in notes if the system records a single primary code.
\end{itemize}
\end{itemize}

Requires patch (change control; revalidation required):

\begin{itemize}
\item Add a new test week (new protected window not already defined in Stage 3.B/3.D).
\item Extend phase duration.
\item Change which tests serve as gate tests for the phase.
\item Change the cadence pattern or work-to-deload ratio for the phase.
\item Change the phase-level stress map structure (protected window placement pattern).
\end{itemize}

Documentation requirements (minimum; applies to all adjustments):

\begin{itemize}
\item Record validity tags (session: \textbf{VALID}/\textbf{MODIFIED}/\textbf{INVALID}; test: \textbf{VALID}/\textbf{INVALID}) and reason-code posture (\textbf{SESSION-RC} or \textbf{TEST-RC}) with free-text context.
\item Record gate impact posture (\textbf{ADVANCE}/\textbf{HOLD}/\textbf{REGRESS}/\textbf{MEDICAL} \textbf{HOLD}) and the reslot target week.
\item When patch is required: create the patch entry per change-control doctrine and record revalidation scope (at minimum Stage 3 artifacts + affected Stage 4–5 phase consumers, plus crosswalk impact check when required).
\end{itemize}

\clearpage
\begin{landscape}

\begingroup
\scriptsize
\setlength{\tabcolsep}{2pt}
\renewcommand{\arraystretch}{1.12}

\begin{longtblr}{
  width=\linewidth,
  rowhead=1,
  colspec = {
    X[c,1.35]
    X[c,1.45]
    Q[c,wd=1.20in]
    X[c,1.25]
    X[c,1.70]
    X[c,2.85]
  },
  hlines,
  vlines,
  cells = {hsep=2pt, vsep=6pt, halign=l, valign=t},
  row{odd} = {bg=white},
  row{even} = {bg=LightGray},
  row{1} = {bg=StageBlue!15, fg=black, font=\bfseries, halign=l, valign=m},
}
\Hdr{Adjustment Type} &
\Hdr{Example} &
\Hdr{Allowed Without Patch} &
\Hdr{Requires Patch} &
\Hdr{Gate Impact} &
\Hdr{Documentation Required} \\

Reslot within phase &
End-of-phase \textbf{RUN} time trial \textbf{INVALID} due to unsafe ice &
Yes &
No &
\textbf{HOLD} for the affected metric until \textbf{VALID} retest; phase remains current &
Test row: \textbf{INVALID} + \textbf{TEST-RC-002} (\textbf{ENV}) or \textbf{TEST-RC-004} (\textbf{SAFETY}); reslot target week; environment notes; Evidence Ref \\

Swap test order inside same protected week &
Move \textbf{CAL} battery earlier in the deload+test week; move \textbf{SWIM} later &
Yes &
No &
None if \textbf{VALID} evidence produced; otherwise \textbf{HOLD} for missing metric &
Updated staged battery note; confirm caps/spacing intent preserved; record any constraints \\

Convert deload+test to deload-only &
Sleep collapse and rising pain trend in the lead-in week &
Yes &
No &
\textbf{HOLD} for unexecuted gate tests; reslot to next protected window &
Weekly review note; sessions tagged appropriately; reason code posture (\textbf{SESSION-RC-003} \textbf{ILL} or \textbf{SESSION-RC-002} PAIN or \textbf{SESSION-RC-007} \textbf{SAFETY} as applicable); reslot target \\

Remove underwater from scheduled window &
Certified lifeguard + dedicated safety spotter not available &
Yes &
No &
Underwater metric remains untested; \textbf{HOLD} for underwater metric; no proxy credit &
Record: “not scheduled; inadmissible without governance; surface-only schedule continues”; do not attempt; no substitution claim \\

Reslot due to travel &
\textbf{TRAVEL} disrupts access to route/pool &
Yes &
No &
\textbf{HOLD} if end gate cannot be executed validly &
Session tags + \textbf{SESSION-RC-006} (\textbf{TRAVEL}) (and \textbf{TRAVEL} note); test reslot target; environment/access notes \\

Add new protected test week &
Insert an extra deload+test week after the planned end gate &
No &
Yes &
Phase extension or schedule change impacts gate timing; default \textbf{HOLD} until executed &
Patch entry; updated phase map; revalidation scope recorded; crosswalk impact check \\

Extend phase duration &
Extend Phase by +2 weeks to regain a protected window &
No &
Yes &
\textbf{HOLD} until standards met with admissible evidence; no \textbf{ADVANCE} until met &
Patch entry; updated Stage 3.B/3.D references; downstream revalidation list \\

Change gate test set &
Replace a \textbf{RUN} gate with an \textbf{AER} time trial as “equivalent” &
No &
Yes (and likely prohibited) &
Prohibited equivalency; cannot justify \textbf{ADVANCE} &
Patch request denied unless Stage 2 authorizes; record non-deployable posture if attempted \\

\end{longtblr}

\endgroup

\end{landscape}
\clearpage

\begin{landscape}

\subsection*{3.E.6 Decision Table}
\phantomsection
\addcontentsline{toc}{subsection}{3.E.6 Decision Table}

\begingroup
\scriptsize
\setlength{\tabcolsep}{2pt}
\renewcommand{\arraystretch}{1.12}

\begin{longtblr}{
  width=\linewidth,
  rowhead=1,
  colspec = {
    X[l,1.65]
    X[l,1.55]
    X[l,1.20]
    X[l,1.85]
    X[l,2.35]
  },
  hlines,
  vlines,
  cells = {hsep=6pt, vsep=6pt, halign=l, valign=t},
  cell{1-Z}{1-5} = {fg=black},
  row{odd} = {bg=white},
  row{even} = {bg=LightGray},
  row{1} = {bg=StageBlue!15, fg=black, font=\bfseries, halign=l, valign=m},
}
\Hdr{Situation} &
\Hdr{Example Trigger} &
\Hdr{Immediate Action} &
\Hdr{Next-Step Adjustment} &
\Hdr{Notes} \\

Minor fatigue with slightly elevated RPE trend &
\textbf{RECOVERY} markers elevated; no red flags &
Proceed with deload posture; keep tests tentative &
Convert deload+test week to check-in-only if trends persist; reslot maximal tests to next protected window &
Do not force maximal attempts under fatigue contamination risk; document \textbf{RECOVERY} trend and scheduling decision; no gate credit without \textbf{VALID} tests \\

Plateau for several weeks with no health concerns &
Multiple \textbf{VALID} results stable but not improving &
Maintain schedule; do not compress &
Keep end-of-phase gate; add only one check-in major event in next protected window if needed for replication &
Plateau does not justify stacking; focus is replication-quality and validity; document as \textbf{HOLD} if standards not met \\

Significant regression on a key \textbf{RUN} or \textbf{RUCK} gate test with \textbf{VALID} result below threshold &
\textbf{VALID} \textbf{RUN}/\textbf{RUCK} result below required tier for current gate &
Record \textbf{VALID} result; do not rationalize upward &
Reslot retest to next protected window; consider deload-only conversion if durability/recovery signals are degraded &
No immediate “retest tomorrow”; preserve spacing intent; gate posture defaults \textbf{HOLD} if end gate unmet \\

Onset of pain with altered mechanics suggesting overuse trend &
\textbf{DURABILITY}: gait-altering pain flag = Yes &
\textbf{STOP} \textbf{NOW} for provoking exposure; terminate test attempt &
Convert week to deload-only; reslot major tests; consider \textbf{REGRESS} posture if mechanics change persists &
Test attempt is \textbf{INVALID} if continued; document pain + mechanics change; sessions may become \textbf{MODIFIED}/\textbf{INVALID} with \textbf{SESSION-RC-002} (PAIN) \\

Illness with fever/systemic symptoms or respiratory illness reducing exertion tolerance &
Systemic illness signs; abnormal dyspnea at low effort &
Cancel major tests; do not “test anyway” &
Convert week to deload-only; reslot to next protected window; default \textbf{HOLD} for missed gate &
Tests aborted are \textbf{INVALID} for gate use; use \textbf{SESSION-RC-003} (\textbf{ILL}) where applicable and \textbf{TEST-RC-004} (\textbf{SAFETY}) for aborted tests \\

Red-flag symptoms on test day (cardiovascular/neurologic/heat illness) &
Chest pain; syncope; confusion; heat illness signs &
\textbf{STOP} \textbf{NOW}; terminate all testing &
\textbf{MEDICAL} \textbf{HOLD} posture when indicated; reslot all gate tests; no make-up stacking &
Any test evidence that continues is inadmissible; document trigger, action, escalation; gate posture \textbf{HOLD}/ \textbf{MEDICAL} \textbf{HOLD} as applicable \\

Environment invalidates testing (route/pool/machine) &
Ice; lightning; poor air quality; pool closure/crowding; machine settings not verifiable &
Do not attempt gate-grade test &
Reslot to next protected window or substitute indoors only if Stage 2 validity/logging preserves comparability &
If attempted under compromised conditions, tag \textbf{INVALID} with \textbf{TEST-RC-002} (\textbf{ENV}) or \textbf{TEST-RC-001} (\textbf{MEAS}); do not claim tier attainment \\

Water governance unavailable for underwater work &
No certified lifeguard and/or no dedicated safety spotter &
Do not attempt underwater; record not authorized &
Remove underwater from window; reslot only when governance is guaranteed &
Document: “not scheduled; inadmissible without governance; surface-only schedule continues.” No proxy equivalency permitted \\

Travel disrupts scheduled test week &
\textbf{TRAVEL} removes access to route/pool/tools &
Reslot tests; preserve safety and validity &
Use next protected window; if end gate missed and no window remains, phase extension under patch control &
Log sessions with \textbf{SESSION-RC-006} (\textbf{TRAVEL}) and \textbf{TRAVEL} note; do not stack tests on return; gate posture defaults \textbf{HOLD} if tests cannot be executed validly \\

\end{longtblr}

\endgroup

\end{landscape}
\clearpage

\subsection*{3.E.7 Conflict Resolution Rule}
\phantomsection
\addcontentsline{toc}{subsection}{3.E.7 Conflict Resolution Rule}

\begin{itemize}
\item Missed, \textbf{INVALID}, or unsafe tests cannot be used for advancement. \textbf{INVALID} provides no gate credit and cannot support \textbf{ADVANCE}.
\item If any adjustment violates Stage 1 governance (validity tags, stop posture, change control) or Stage 2 validity requirements (protocol conditions, logging completeness), downstream execution is Draft — Non-Deployable until corrected and re-versioned.
\item Safety overrides schedule and performance. \textbf{STOP} \textbf{NOW} and \textbf{MEDICAL} \textbf{HOLD} postures control scheduling decisions when triggers occur.
\end{itemize}

\end{document}